\subsection{Bulk Neutrinos as an Alternative Cause of the Gallium and Reactor Anti-neutrino Anomalies --- arXiv:1107.2400}

\textbf{Authors:} P. A. N. Machado, H. Nunokawa, F. A. Pereira dos Santos, R. Zukanovich Funchal

\paragraph{主な主張}
Gallium線源実験と19の短基線原子炉実験の電子ニュートリノ・反ニュートリノ欠損をKK sterile状態への振動で説明でき、当時の他の振動制約とも大部分で両立する領域があると示す\cite{Machado:2012LEDAnomalies}。

\paragraph{新規性}
独立な質量・混合を持つ少数のsterile状態を追加する代わりに、ニュートリノ質量生成と結び付いたKK塔で二種類のdisappearance anomalyを共通に説明した。

\paragraph{位置付け}
LEDによるGallium・reactor anomaly解釈の代表的な初期提案であり、その適合性は後年の実験データによる再評価と区別する必要がある。
